\documentclass[../../../include/open-logic-section]{subfiles}

\begin{document}

\olfileid{sth}{replacement}{extrinsic}
\olsection[బాహ్య పరిశీలనలు]{ప్రతిస్థాపన గురించిన బాహ్య పరిశీలనలు}

మొదట ప్రతిస్థాపనను సమర్థించే \emph{బాహ్య} ప్రయత్నాన్ని పరిశీలిద్దాం. బూలాస్ ఒకదాన్ని ఇలా సూచిస్తారు.
\begin{quote}
  [\ldots] ప్రతిస్థాపన స్వయంసిద్ధ సూత్రాలను స్వీకరించడానికి కారణం చాలా సరళం: వాటికి ఎన్నో కోరదగిన ఫలితాలు ఉన్నాయి; (కనిపించినంతవరకు) కోరనివి లేవు. పునరావర్తన భావన గురించిన సిద్ధాంతాలతో పాటు, వాటి ఫలితాల్లో అనంత సంఖ్యలపై ఆదర్శవంతమైనది కాకపోయినా సంతృప్తికరమైన సిద్ధాంతం, సుపునాది సంబంధాలపై ఆగమన నిర్వచనాలను సమర్థించే ఎంతో కోరదగిన ఫలితం ఉన్నాయి.
  \citep[229]{Boolos1971}
\end{quote}
బూలాస్ ఆలోచన సారాంశం: ప్రతిస్థాపనను దాని ఫలితాలతో సమర్థించాలి. ఆయన ప్రస్తావించిన నిర్దిష్ట ఫలితాలు గత కొన్ని అధ్యాయాల్లో మనం చర్చించినవే. వాన్ నాయ్‌మన్ క్రమసంఖ్యలు సుక్రమరకం భావనకు చక్కని ప్రతినిధులని ప్రతిస్థాపనతో నిరూపించాం (ఇదే మన ``అనంత సంఖ్యలపై ఆదర్శవంతమైనది కాకపోయినా సంతృప్తికరమైన సిద్ధాంతం''). ప్రతిస్థాపనతోనే $V_\alpha$లను నిర్వచించి, స్థాయి భావనను స్థాపించి, $\in$-ఆగమనాన్ని నిరూపించాం (ఇవే మన ``పునరావర్తన భావన గురించిన సిద్ధాంతాలు''). చివరగా, ప్రతిస్థాపనతో అతిపరిమిత పునరావృత్తి సిద్ధాంతాన్ని నిరూపించవచ్చు (ఇదే ``సుపునాది సంబంధాలపై ఆగమన నిర్వచనాలు'').

ఇవి నిజంగానే కోరదగిన ఫలితాలు. కానీ ప్రతిస్థాపనను \emph{సమర్థించడానికి} ఇవి సరిపోతాయా? \emph{కాదు}. కనీసం నేరుగా సరిపోవు.

ఒక సరళమైన సమస్య ఇది. ప్రతిస్థాపనకు కొన్ని కోరదగిన ఫలితాలు ఉన్నాయని చెప్పినా, వాటిలో ఎన్నో ఇతర మార్గాల్లో పొందవచ్చు. ఇది తెలిసినంతగా తెలియని విషయం; కనుక పరిస్థితిని కాసేపు వివరించాలి.

స్థర సిద్ధాంతం, సంక్షిప్తంగా $\LT$, అనే సరళమైన సమితి సిద్ధాంతం ఉంది.\footnote{$\LT$ తొలి రూపాలను \citet{Montague1965}, \citet{Scott1974} ఇచ్చారు; \citet{Potter2004} దాన్ని సరళీకరించి పుస్తకస్థాయిలో చర్చించారు; ఇటీవల \citet{ButtonLT1} $\LT$ను ఇంకా సరళీకరించారు.} $\LT$ స్వయంసిద్ధ సూత్రాలు విస్తరణీయత, వేరుచేయడం, ప్రతి సమితీ ఏదో ఒక \emph{స్థరానికి} ఉపసమితి అనే వాదన మాత్రమే. ఇక్కడ ``స్థరం''ను చాకచక్యంగా నిర్వచిస్తారు; ఆ స్థరాలు మన $V_\alpha$ల్లా ప్రవర్తిస్తాయి. అందువల్ల $\ZF$తో $\LT$ను నిరూపించవచ్చు; అయితే $\LT$, $\ZF$కంటే \emph{చాలా} బలహీనమైనది. నిజానికి $\LT$ జతలు, ఘాత సమితులు, అనంతత్వం, ప్రతిస్థాపన ఏదీ ఇవ్వదు. $\LT$కు అనంతత్వం, ఘాత సమితులను చేర్చగా వచ్చే సిద్ధాంతాన్ని $\Zr$ అనుకుందాం; దీని నుండి జతలు కూడా వస్తాయి, కాబట్టి $\Zr$ కనీసం $\Z$ంత బలమైనది. వాస్తవానికి, ప్రతి సమితికీ స్థాయి ఉందనే వాదనను చేర్చడం వల్ల $\Zr$, $\Z$కంటే ఖచ్చితంగా బలమైనది (దాన్ని $\Zr$ అని పిలవాలనే నా సూచనకు ఇదే కారణం). నిజానికి $\Zr$ ఇచ్చేవి: క్రమసంఖ్యలపై పూర్తిగా సంతృప్తికరమైన సిద్ధాంతం;
స్థాయిక్రమాన్ని సుక్రమబద్ధ దశలుగా విభజించే ఫలితాలు; $\in$-ఆగమన నిరూపణ; అతిపరిమిత పునరావృత్తి యొక్క ఒక \emph{రూపం}.

సంక్షిప్తంగా: బూలాస్‌కు ఇది తెలియకపోయినా, ఆయన చెప్పిన కోరదగిన ఫలితాలన్నీ ప్రతిస్థాపన \emph{లేకుండానే} పొందవచ్చు; $\Z$ బదులు $\Zr$ వాడితే సరిపోతుంది.

(ఇదంతా తెలిసి కూడా, $\LT$, $\Zr$ బదులు సంప్రదాయ మార్గంలో $\ZF$నే మీకు ఎందుకు బోధించాం? రెండు కారణాలు. మొదట: పూర్తిగా చారిత్రక కారణాల వల్ల $\LT$తో మొదలుపెట్టడం సాధారణ పద్ధతి కాదు; సమితి సిద్ధాంతపు ప్రామాణిక చర్చలను మీరు చదవగలిగేలా చేయాలనుకున్నాం. రెండవది: $\LT$, $\Zr$ను గ్రహించడానికి మీరు సిద్ధమైనప్పుడు \citealt{Potter2004}, \citealt{ButtonLT1}ను చదవవచ్చు.)

నిజమే, $\Zr$, $\ZF$కంటే ఖచ్చితంగా బలహీనమైనది కనుక, $\ZF$ నిరూపించే కానీ $\Zr$ నిర్ణయించని ఫలితాలు ఉన్నాయి. వాటిలో ఒకదాన్ని చూపి బాహ్య ఆధారాలపై ప్రతిస్థాపనను సమర్థించవచ్చు. కానీ $\Zr$ను వాడడం నేర్చుకున్నాక, (a) ఇతరంగా కాక ప్రతిస్థాపనతోనే నిర్ణయించబడే, (b) సహజంగా నిజమనిపించే విషయాల ఉదాహరణలు ఎన్నో కనుగొనడం కష్టం. (దీనిపై మరింతకు \citealt[\S13.2]{Potter2004} చూడండి.)

చివరి విషయం ఇది: ప్రతిస్థాపనకు బలమైన బాహ్య సమర్థన ఇవ్వాలంటే, ప్రతిస్థాపన \emph{లేకుండా} సాధించలేని ఫలితాన్ని కనుగొనాలి. అది సులభమైన పని కాదు.

పూర్తిగా బాహ్య సమర్థన ప్రయత్నానికి ఎదురయ్యే మరో సమస్యను పరిశీలిద్దాం. (ఇది మొదటిదానికంటే ప్రాథమికమైనదీ కావచ్చు.) ప్రతిస్థాపనకు ఎన్నో కోరదగిన ఫలితాలున్నాయని మాత్రమే బూలాస్ చెప్పరు. దానికి ``(కనిపించినంతవరకు) కోరని'' ఫలితాలు లేవని కూడా అంటారు. అయితే కుండలీకరణంలోని ``కనిపించినంతవరకు'' అనే పరిమితి ఖచ్చితంగా అత్యంత ముఖ్యమైనది.

మనం ఇక్కడికి ఎలా వచ్చామో గుర్తుచేసుకోండి: అమాయక అవగాహన సూత్రం అసంగతికి దారితీసింది; దానికి ప్రతిస్పందనగా సమితి యొక్క సంచిత-పునరావర్తన భావనను స్వీకరించాం. ఈ భావనతో కూడిన కథనం, అది సుసంగతమని మనకు భరోసా ఇస్తుందని ఆశిస్తున్నాం. కానీ ఆ కథనంలోనే ప్రతిస్థాపనను సమర్థించలేకపోతే, $\ZF$ సుసంగతమని నమ్మడానికి (ఇంతవరకు) కారణం లేదు. లేదా మరింత ఖచ్చితంగా: \emph{ఇప్పటివరకు} ఎవరూ వైరుధ్యాన్ని కనుగొనలేదనే (బహుశా యాదృచ్ఛిక) విషయం తప్ప $\ZF$ సుసంగతమని నమ్మడానికి కారణం లేదు. ఆ అర్థంలో బూలాస్ వ్యాఖ్య చివరకు ``(కనిపించినంతవరకు) $\ZF$ సుసంగతమే'' అన్నట్టవుతుంది. సుసంగతత గురించి దీనికంటే బలమైన భరోసా కోరాలి.

ఈ సమస్య ప్రతిస్థాపనకు \emph{పూర్తిగా} బాహ్య సమర్థన ఇచ్చే ఏ ప్రయత్నానికైనా వర్తిస్తుంది; అంటే $\ZF$ యొక్క (తెలిసిన) ఫలితాల ఆధారంగానే చెప్పే సమర్థనకు. కాబట్టి ప్రతిస్థాపనకు \emph{అంతర్గత} సమర్థనను వెతకాలి; అంటే సమితుల గురించి మనం చెప్పిన కథనం ప్రతిస్థాపనకు ఇప్పటికే ఏదోలా మనలను బద్ధులను చేస్తుందని సూచించే సమర్థనను.

\end{document}
